\chapter{The Loss Landscape of RL}
\label{ch:loss-landscape}

Reinforcement learning borrows the optimization machinery of supervised learning---forward pass, loss computation, backpropagation, parameter update---but the losses it optimizes are fundamentally different. In supervised learning, the loss measures deviation from externally provided labels on a fixed dataset. In RL, the ``labels'' are constructed by the agent itself, the data distribution shifts as the policy changes, and the loss must encode not just fitting accuracy but also exploration incentives and distributional control.

This chapter traces the evolution of RL loss functions: from MSE regression on Bellman targets, through direct Q-maximization, to weighted log-likelihood and entropy regularization. Each transition addresses a specific failure mode of the previous formulation.

%% ----------------------------------------------------------------
\section{Value-Based Loss: MSE on Self-Generated Labels}
\label{sec:td-mse}

The starting point is deceptively simple. DQN and Q-learning frame RL as a regression problem:
\begin{equation}
  \mathcal{L}_{\text{critic}} = \bigl(Q_\theta(s,a) - y\bigr)^2,
  \qquad
  y = r + \gamma \max_{a'} Q_{\bar\theta}(s', a'),
  \label{eq:td-loss}
\end{equation}
where $\bar\theta$ denotes the parameters of a slowly updated target network.

\begin{keyidea}[Self-Bootstrapping Regression]
The target $y$ is not an external ground-truth label. It is computed by the model itself through the Bellman bootstrap and the $\max$ operator. The model creates its own labels, then fits those labels. This is the core characteristic of TD learning and the reason it is \emph{not} ordinary MSE.
\end{keyidea}

Two consequences follow immediately:
\begin{enumerate}
  \item \textbf{The target is noisy.} Because $Q_{\bar\theta}$ is an approximation, $y$ carries estimation error that feeds back into training.
  \item \textbf{The $\max$ operator amplifies noise.} Among noisy estimates, the maximum is biased upward (Section~\ref{sec:overestimation}).
\end{enumerate}

%% ----------------------------------------------------------------
\section{Q Overestimation}
\label{sec:overestimation}

Suppose each action's Q-estimate decomposes as
\begin{equation}
  \hat{Q}(a) = Q^*(a) + \epsilon(a),
  \qquad \epsilon(a) \sim \text{noise}.
\end{equation}
When we take the maximum over actions, the action whose noise $\epsilon(a)$ happens to be largest is preferentially selected. By Jensen's inequality applied to the convex $\max$ function:
\begin{equation}
  \mathbb{E}\!\bigl[\max_a \hat{Q}(a)\bigr]
  \;\ge\;
  \max_a \mathbb{E}\!\bigl[\hat{Q}(a)\bigr].
  \label{eq:overestimation}
\end{equation}

This is not a peculiarity of DQN. \emph{Any} algorithm that relies on $\max Q$, $\arg\max Q$, or advantage ranking faces the same issue. The difference between algorithms lies in how aggressively they mitigate it.

\begin{keyidea}[Overestimation Is Structural]
Q overestimation is not a bug in a particular implementation. It is a mathematical consequence of taking the maximum of noisy estimates. The more actions, the larger the bias.
\end{keyidea}

%% ----------------------------------------------------------------
\section{Double DQN and TD3: Stabilizing the Target}
\label{sec:double-td3}

\subsection{Double DQN}

The core idea is to \emph{decouple action selection from value evaluation}. Standard DQN uses the same (noisy) network for both:
\begin{equation}
  y_{\text{DQN}} = r + \gamma \max_{a'} Q_{\bar\theta}(s', a').
\end{equation}
Double DQN selects the action with the online network but evaluates it with the target network:
\begin{equation}
  a^* = \arg\max_{a'} Q_\theta(s', a'),
  \qquad
  y_{\text{DDQN}} = r + \gamma\, Q_{\bar\theta}(s', a^*).
  \label{eq:double-dqn}
\end{equation}
This does not eliminate bias entirely, but it breaks the positive feedback loop in which the same noisy Q is responsible for both choosing and scoring.

\subsection{TD3: Clipped Double-Q for Continuous Actions}

TD3 extends the idea to continuous action spaces. It maintains two independent critics $Q_{\phi_1}$, $Q_{\phi_2}$ and uses the more pessimistic estimate:
\begin{equation}
  y = r + \gamma \min\bigl(Q_{\phi_1}(s', \tilde{a}'),\; Q_{\phi_2}(s', \tilde{a}')\bigr),
  \qquad
  \tilde{a}' = \mu_{\bar\theta}(s') + \text{clip}(\epsilon, -c, c).
  \label{eq:td3-target}
\end{equation}
The $\min$ produces a conservative estimate that counteracts the actor's tendency to exploit critic errors. Two additional mechanisms---\emph{delayed policy updates} (let the critic stabilize before updating the actor) and \emph{target policy smoothing} (add clipped noise to the target action)---further prevent the actor from chasing sharp, unreliable peaks in the Q landscape.

%% ----------------------------------------------------------------
\section{Max-Q Loss: The Actor Directly Maximizes Q}
\label{sec:max-q}

In discrete action spaces, $\arg\max_a Q(s,a)$ is trivial: enumerate all actions, pick the winner. In continuous spaces this is intractable, so DDPG introduces an actor network $\mu_\theta(s)$ trained to output the action that maximizes the critic:
\begin{equation}
  \mathcal{L}_{\text{actor}} = -\,Q_\phi\!\bigl(s,\, \mu_\theta(s)\bigr).
  \label{eq:ddpg-actor}
\end{equation}
Minimizing this loss is equivalent to gradient ascent on $Q$ with respect to the actor's output. The gradient signal flows through the critic: $\nabla_\theta \mathcal{L}_{\text{actor}} = -\nabla_a Q_\phi(s,a)\big|_{a=\mu_\theta(s)} \cdot \nabla_\theta \mu_\theta(s)$.

\begin{keyidea}[Differentiable Argmax]
DDPG's actor loss turns the discrete, non-differentiable $\arg\max$ operation into a differentiable gradient optimization problem. The actor uses the critic's gradient landscape $\nabla_a Q$ to steer its output toward high-value actions.
\end{keyidea}

\subsection{The Fundamental Risk}

The critic's Q is not the true value function---it has noise, bias, and overestimation. If the actor relentlessly climbs the Q surface, it may converge not to truly good actions but to actions where the critic's \emph{estimation error} is largest. This is why TD3 requires its suite of protective measures (double critic, delayed updates, target smoothing).

\subsection{Placing Max-Q in the Loss Evolution}

\begin{center}
\begin{tabular}{@{} l l @{}}
\toprule
\textbf{Loss type} & \textbf{What it does} \\
\midrule
MSE (TD) & Make $Q$ approximate a Bellman target --- passive fitting \\
Max-Q (DDPG actor) & Make the actor output high-$Q$ actions --- actively exploiting the Q landscape \\
Policy gradient & Make good actions more probable --- directly shaping the distribution \\
\bottomrule
\end{tabular}
\end{center}

MSE says ``I estimate value.'' Max-Q says ``I exploit my estimates for decisions.'' Policy gradient says ``I directly optimize decisions.''

%% ----------------------------------------------------------------
\section{Policy Gradient: Gradient First, Loss Second}
\label{sec:pg-loss}

Policy gradient methods bypass the Q landscape entirely. The objective is the expected return:
\begin{equation}
  J(\theta) = \mathbb{E}_{\pi_\theta}\!\bigl[\text{return}\bigr].
\end{equation}
The policy gradient theorem yields:
\begin{equation}
  \nabla_\theta J(\theta)
  = \mathbb{E}_{s,a\sim\pi_\theta}\!\bigl[
    \nabla_\theta \log\pi_\theta(a\mid s)\;\hat{A}(s,a)
  \bigr],
  \label{eq:pg}
\end{equation}
where $\hat{A}(s,a)$ is an advantage estimate.

\subsection{Why ``Gradient First''}

A common source of confusion: value-based methods start from a loss and derive a gradient, but policy gradient derivations start from the gradient and only mention a loss at the end. The reason is structural.

\textbf{Value-based:} There is a natural regression target $y$. The loss $\mathcal{L} = (Q - y)^2$ is the starting point; the gradient follows from differentiation.

\textbf{Policy gradient:} The starting point is the objective $J(\theta)$. The gradient~\eqref{eq:pg} is derived via the likelihood ratio trick applied to trajectory probabilities. No loss function is written first.

To use automatic differentiation frameworks, we \emph{construct a surrogate loss after the fact}:
\begin{equation}
  \mathcal{L}_{\text{PG}} = -\log\pi_\theta(a\mid s)\;\hat{A}(s,a).
  \label{eq:pg-loss}
\end{equation}
Differentiating this with respect to $\theta$ recovers exactly the policy gradient~\eqref{eq:pg} (with the sign flip because optimizers minimize). The loss is an engineering artifact for autograd, not the conceptual starting point.

\begin{keyidea}[Constructed Loss]
In value-based RL, loss comes first and gradient follows. In policy gradient RL, the gradient comes first from mathematical derivation; the loss is constructed afterward so that \texttt{loss.backward()} produces the correct gradient.
\end{keyidea}

Yet in code, both families share the same training loop: \texttt{loss = compute\_loss(batch); loss.backward(); optimizer.step()}.

%% ----------------------------------------------------------------
\section{Connection to Maximum Likelihood}
\label{sec:mle-connection}

\subsection{Behavior Cloning as Unweighted MLE}

Given demonstration data $\{(s_i, a_i)\}$, behavior cloning maximizes:
\begin{equation}
  \sum_i \log\pi_\theta(a_i \mid s_i).
\end{equation}
Every observed action receives equal weight. The model learns ``a human did this,'' with no notion of whether ``this'' was good or bad.

\subsection{Policy Gradient as Weighted MLE}

Policy gradient adds a quality signal:
\begin{equation}
  \mathcal{L}_{\text{PG}} = -\hat{A}(s,a)\;\log\pi_\theta(a \mid s).
\end{equation}
The advantage $\hat{A}$ serves as a per-sample weight. Actions that led to higher returns get amplified; actions that underperformed get suppressed.

\begin{keyidea}[From MLE to Weighted MLE]
Behavior cloning (unweighted MLE) models behavioral frequency: ``this action appeared, so increase its probability.'' Policy gradient (weighted MLE) models behavioral frequency with preference weights: ``this action appeared \emph{and} brought high returns, so increase its probability proportionally.''
\end{keyidea}

This perspective makes the conceptual shift transparent:
\begin{center}
\begin{tabular}{@{} l l @{}}
\toprule
\textbf{Method} & \textbf{What it models} \\
\midrule
MLE / Behavior Cloning & Frequency of observed actions \\
Policy Gradient & Frequency weighted by quality of outcomes \\
\bottomrule
\end{tabular}
\end{center}

%% ----------------------------------------------------------------
\section{Entropy: Sampling Distribution Regularization}
\label{sec:entropy}

PPO and SAC both include an entropy bonus $\mathcal{H}[\pi_\theta(\cdot \mid s)]$. Its role is often trivialized as ``encouraging randomness.'' The real function is more precise: it prevents the policy distribution from collapsing prematurely.

\subsection{The Collapse Mechanism}

Without entropy regularization, the policy gradient pushes probability mass toward actions with high estimated advantage. The resulting feedback loop is:
\begin{equation*}
  \text{policy concentrates}
  \;\to\; \text{sampling narrows}
  \;\to\; \text{data diversity drops}
  \;\to\; \text{Q/advantage estimates degrade}
  \;\to\; \text{local optima.}
\end{equation*}

\subsection{Entropy as a Distributional Force}

Policy optimization involves two competing forces:

\begin{center}
\begin{tabular}{@{} l l @{}}
\toprule
\textbf{Force} & \textbf{Effect} \\
\midrule
Preference term $\hat{A}\cdot\log\pi$ & Push probability toward high-advantage actions \\
Entropy term $\mathcal{H}[\pi]$ & Pull the distribution toward uniformity \\
\bottomrule
\end{tabular}
\end{center}

The learned policy is the equilibrium between these two forces---neither a uniform distribution nor a collapsed delta, but a controlled trade-off.

\subsection{SAC: Entropy as Part of the Objective}

SAC goes further by writing the objective as:
\begin{equation}
  J_{\text{SAC}}(\theta) = \mathbb{E}\!\Bigl[\sum_t r_t + \alpha\,\mathcal{H}\bigl[\pi_\theta(\cdot \mid s_t)\bigr]\Bigr],
  \label{eq:sac-obj}
\end{equation}
where $\alpha$ is a temperature parameter (often auto-tuned). Exploration is no longer an external trick appended to the loss; it is a first-class component of the optimization objective. Sampling diversity is something the agent is \emph{rewarded} for.

%% ----------------------------------------------------------------
\section{The Full Loss Evolution}
\label{sec:loss-evolution}

We can now see the complete arc:

\begin{center}
\begin{tabular}{@{} l l l @{}}
\toprule
\textbf{Loss layer} & \textbf{Formula sketch} & \textbf{What it controls} \\
\midrule
MSE (TD) & $(Q - y)^2$ & Value estimation accuracy \\
Max-Q (actor) & $-Q(s, \mu_\theta(s))$ & Action quality via critic landscape \\
Weighted $\log\pi$ & $-\hat{A}\cdot\log\pi_\theta(a \mid s)$ & Preference-weighted behavior \\
Entropy & $+\alpha\,\mathcal{H}[\pi]$ & Sampling diversity control \\
\bottomrule
\end{tabular}
\end{center}

These four layers correspond to four progressively deeper concerns:
\begin{enumerate}
  \item \textbf{Estimate value} (how much is this worth?).
  \item \textbf{Exploit estimates} (given the value landscape, what should I do?).
  \item \textbf{Shape the policy} (make good actions more probable).
  \item \textbf{Control the distribution} (keep the policy from collapsing).
\end{enumerate}

Modern algorithms like SAC combine all four layers---a critic trained with TD-MSE, an actor that maximizes entropy-augmented Q, and an entropy term that regularizes the policy distribution---into a single coherent framework. Understanding each layer in isolation makes the composite loss transparent.
